\section{Circuit Complexity --- Foundations}
%%% generated by the blueprint pipeline from TCSlib.Complexity.CircuitComplexity.Basic
\subsection*{Circuit Complexity --- Foundations: Introduction}

Declarations of the Lean module \texttt{TCSlib.Complexity.CircuitComplexity.Basic}.

\subsection*{Circuit Complexity --- Foundations: Results}

\begin{theorem}[Evaluation of a literal circuit]
\lean{BoolCircuit.TreeCircuit.eval_lit}
\label{BoolCircuit.TreeCircuit.eval_lit}
\leanok
\uses{BoolCircuit.Lit, BoolCircuit.Lit.eval, BoolCircuit.TreeCircuit.eval}
\difficulty{1}
For every natural number $n$, every literal $l$ on $n$ variables, and every Boolean
assignment $x$ to those variables, the circuit consisting of the single leaf $l$
evaluates to the value of $l$ under $x$.
\end{theorem}

\begin{theorem}[True evaluation of an AND gate]
\lean{BoolCircuit.TreeCircuit.eval_node_true_iff}
\label{BoolCircuit.TreeCircuit.eval_node_true_iff}
\leanok
\uses{BoolCircuit.TreeCircuit, BoolCircuit.TreeCircuit.eval}
\difficulty{4}
For any natural number $n$, any finite list $cs$ of Boolean circuit trees on $n$
variables, and any assignment $x : \mathrm{Fin}\,n \to
\{\mathrm{false},\mathrm{true}\}$, the AND gate with children $cs$ evaluates to true
under $x$ if and only if every circuit in $cs$ evaluates to true under $x$.
\end{theorem}

\begin{theorem}[True evaluation of an OR node]
\lean{BoolCircuit.TreeCircuit.eval_node_false_iff}
\label{BoolCircuit.TreeCircuit.eval_node_false_iff}
\leanok
\uses{BoolCircuit.TreeCircuit, BoolCircuit.TreeCircuit.eval}
\difficulty{4}
For any natural number $n$, list $cs$ of circuits on $n$ variables, and assignment $x$
to those variables, the OR node with children $cs$ evaluates to true if and only if
there exists a circuit in $cs$ that evaluates to true under $x$.
\end{theorem}

\begin{definition}[Maximum fan-in of a list of circuits]
\lean{BoolCircuit.TreeCircuit.maxFaninL}
\statementsource{arora-barak-ch13-circuit-lowerbounds}{pdf-52929dbaa962-p001-b003}
\label{BoolCircuit.TreeCircuit.maxFaninL}
\leanok
\uses{BoolCircuit.TreeCircuit, BoolCircuit.TreeCircuit.maxFanin}
For a list of Boolean circuit trees, define its maximum fan-in to be the largest maximum
fan-in of any circuit in the list, taking the value to be $0$ for the empty list.
\end{definition}

\begin{theorem}[Depth of a circuit node]
\lean{BoolCircuit.TreeCircuit.depth_node}
\label{BoolCircuit.TreeCircuit.depth_node}
\leanok
\uses{BoolCircuit.TreeCircuit, BoolCircuit.TreeCircuit.depth, BoolCircuit.TreeCircuit.maxDepth}
\difficulty{1}
For any natural number $n$, Boolean value $b$, and list $cs$ of circuits on $n$
variables, the depth of the node with gate type $b$ and children $cs$ is one plus the
maximum depth among the children, where the maximum is zero when the list is empty.
\end{theorem}

\begin{theorem}[Size of a gate circuit]
\lean{BoolCircuit.TreeCircuit.size_node}
\label{BoolCircuit.TreeCircuit.size_node}
\leanok
\uses{BoolCircuit.TreeCircuit, BoolCircuit.TreeCircuit.size, BoolCircuit.TreeCircuit.sumSize}
\difficulty{1}
For every natural number $n$, every Boolean gate type (AND or OR), and every list $cs$
of circuits on $n$ variables, the circuit with that gate and children $cs$ has size $1$
plus the sum of the sizes of the circuits in $cs$.
\end{theorem}

\begin{theorem}[Max fanin node]
\lean{BoolCircuit.TreeCircuit.maxFanin_node}
\statementsource{arora-barak-ch13-circuit-lowerbounds}{pdf-52929dbaa962-p001-b003}
\label{BoolCircuit.TreeCircuit.maxFanin_node}
\leanok
\uses{BoolCircuit.TreeCircuit, BoolCircuit.TreeCircuit.maxFanin, BoolCircuit.TreeCircuit.maxFaninL}
\difficulty{1}
% dataset-pipeline: stub -- prose pending build_informal_statements.py
Statement pending regeneration.
\end{theorem}

\begin{theorem}[Maximum depth of the empty list]
\lean{BoolCircuit.TreeCircuit.maxDepth_nil}
\label{BoolCircuit.TreeCircuit.maxDepth_nil}
\leanok
\uses{BoolCircuit.TreeCircuit, BoolCircuit.TreeCircuit.maxDepth}
\difficulty{1}
The maximum depth among an empty list of Boolean circuit trees on $n$ variables is $0$.
\end{theorem}

\begin{theorem}[Maximum depth of a circuit list]
\lean{BoolCircuit.TreeCircuit.maxDepth_cons}
\label{BoolCircuit.TreeCircuit.maxDepth_cons}
\leanok
\uses{BoolCircuit.TreeCircuit, BoolCircuit.TreeCircuit.depth, BoolCircuit.TreeCircuit.maxDepth}
\difficulty{1}
Let $n$ be a natural number, let $c$ be a Boolean circuit tree on $n$ variables, and let
$cs$ be a finite list of such trees. The maximum depth among the trees in the list
formed by placing $c$ before $cs$ is the maximum of the depth of $c$ and the maximum
depth among the trees in $cs$, with the maximum depth of an empty list taken to be $0$.
\end{theorem}

\begin{theorem}[Size sum of the empty list]
\lean{BoolCircuit.TreeCircuit.sumSize_nil}
\label{BoolCircuit.TreeCircuit.sumSize_nil}
\leanok
\uses{BoolCircuit.TreeCircuit, BoolCircuit.TreeCircuit.sumSize}
\difficulty{1}
The sum of the sizes in the empty list of Boolean circuits on $n$ variables is $0$.
\end{theorem}

\begin{theorem}[Size sum over a cons list]
\lean{BoolCircuit.TreeCircuit.sumSize_cons}
\label{BoolCircuit.TreeCircuit.sumSize_cons}
\leanok
\uses{BoolCircuit.TreeCircuit, BoolCircuit.TreeCircuit.size, BoolCircuit.TreeCircuit.sumSize}
\difficulty{1}
For a Boolean circuit $c$ on $n$ variables and a finite list $cs$ of such circuits, the
sum of the sizes in the list $c::cs$ equals $\mathrm{size}(c)$ plus the sum of the sizes
in $cs$.
\end{theorem}

\begin{theorem}[Empty-list maximum fanin]
\lean{BoolCircuit.TreeCircuit.maxFaninL_nil}
\label{BoolCircuit.TreeCircuit.maxFaninL_nil}
\leanok
\uses{BoolCircuit.TreeCircuit, BoolCircuit.TreeCircuit.maxFaninL}
\difficulty{1}
For any number $n$ of input variables, the maximum fanin aggregated over the empty list
of Boolean circuits on $n$ variables is $0$.
\end{theorem}

\begin{theorem}[Maximum fanin of a list with a new first circuit]
\lean{BoolCircuit.TreeCircuit.maxFaninL_cons}
\label{BoolCircuit.TreeCircuit.maxFaninL_cons}
\leanok
\uses{BoolCircuit.TreeCircuit, BoolCircuit.TreeCircuit.maxFanin, BoolCircuit.TreeCircuit.maxFaninL}
\difficulty{1}
For every natural number $n$, every Boolean circuit tree $c$ on $n$ variables, and every
finite list $cs$ of Boolean circuit trees on $n$ variables, the maximum fanin of the
list formed by placing $c$ first in $cs$ is the maximum of the fanin of $c$ and the
maximum fanin in $cs$.
\end{theorem}

\begin{theorem}[Depth bounded by maximum list depth]
\lean{BoolCircuit.TreeCircuit.depth_le_maxDepth}
\label{BoolCircuit.TreeCircuit.depth_le_maxDepth}
\leanok
\uses{BoolCircuit.TreeCircuit, BoolCircuit.TreeCircuit.depth, BoolCircuit.TreeCircuit.maxDepth}
\difficulty{4}
For any natural number $n$, any Boolean circuit $c$ on $n$ variables, and any list $cs$
of such circuits containing $c$, the depth of $c$ is at most the maximum depth among the
circuits in $cs$, with the maximum taken to be zero for the empty list.
\end{theorem}

\begin{theorem}[Circuit fanin bounded by list maximum]
\lean{BoolCircuit.TreeCircuit.maxFanin_le_maxFaninL}
\label{BoolCircuit.TreeCircuit.maxFanin_le_maxFaninL}
\leanok
\uses{BoolCircuit.TreeCircuit, BoolCircuit.TreeCircuit.maxFanin, BoolCircuit.TreeCircuit.maxFaninL}
\difficulty{4}
Let $n$ be a natural number, let $c$ be a Boolean circuit on $n$ variables, and let $cs$
be a list of such circuits. If $c$ occurs in $cs$, then the maximum fanin of $c$ is at
most the maximum fanin aggregated over $cs$.
\end{theorem}

\begin{theorem}[List length bounded by total circuit size]
\lean{BoolCircuit.TreeCircuit.length_le_sumSize}
\label{BoolCircuit.TreeCircuit.length_le_sumSize}
\leanok
\uses{BoolCircuit.TreeCircuit, BoolCircuit.TreeCircuit.one_le_size, BoolCircuit.TreeCircuit.sumSize, BoolCircuit.TreeCircuit.sumSize_cons}
\difficulty{4}
For every finite list $cs$ of Boolean circuit trees on $n$ variables, the length of $cs$
is at most the sum of the sizes of its circuits.
\end{theorem}

\begin{theorem}[List maximum fanin bounded by total size]
\lean{BoolCircuit.TreeCircuit.maxFaninL_le_sumSize}
\label{BoolCircuit.TreeCircuit.maxFaninL_le_sumSize}
\leanok
\uses{BoolCircuit.TreeCircuit, BoolCircuit.TreeCircuit.maxFanin, BoolCircuit.TreeCircuit.maxFaninL, BoolCircuit.TreeCircuit.maxFaninL_cons, BoolCircuit.TreeCircuit.size, BoolCircuit.TreeCircuit.sumSize, BoolCircuit.TreeCircuit.sumSize_cons}
\difficulty{4}
Let $cs$ be a finite list of Boolean circuit trees on $n$ variables, and suppose each
$c$ in the list has maximum fanin at most its size. Then the maximum fanin among the
list is at most the sum of the sizes of its circuits.
\end{theorem}

\begin{theorem}[Fanin bounded by circuit size]
\lean{BoolCircuit.TreeCircuit.maxFanin_le_size}
\label{BoolCircuit.TreeCircuit.maxFanin_le_size}
\leanok
\uses{BoolCircuit.TreeCircuit, BoolCircuit.TreeCircuit.ind, BoolCircuit.TreeCircuit.length_le_sumSize, BoolCircuit.TreeCircuit.maxFanin, BoolCircuit.TreeCircuit.maxFaninL_le_sumSize, BoolCircuit.TreeCircuit.maxFanin_node, BoolCircuit.TreeCircuit.size, BoolCircuit.TreeCircuit.size_node}
\difficulty{4}
For every natural number $n$ and every Boolean circuit $c$ on $n$ variables, the maximum
fanin of $c$ is at most its size.
\end{theorem}

\begin{theorem}[Bound on the total size of a circuit list]
\lean{BoolCircuit.TreeCircuit.sumSize_add_length_le}
\label{BoolCircuit.TreeCircuit.sumSize_add_length_le}
\leanok
\uses{BoolCircuit.TreeCircuit, BoolCircuit.TreeCircuit.size, BoolCircuit.TreeCircuit.sumSize, BoolCircuit.TreeCircuit.sumSize_cons, BoolCircuit.TreeCircuit.sumSize_nil}
\difficulty{4}
Let $m$ be a natural number and let $cs$ be a finite list of Boolean circuits on $n$
variables. If each circuit $c$ in $cs$ satisfies $\mathrm{size}(c)+1\le m$, then the sum
of their sizes plus the length of the list is at most the length of the list multiplied
by $m$.
\end{theorem}

\begin{theorem}[Size bound for bounded-fanin circuits]
\lean{BoolCircuit.TreeCircuit.size_succ_le_two_pow}
\label{BoolCircuit.TreeCircuit.size_succ_le_two_pow}
\leanok
\uses{BoolCircuit.TreeCircuit, BoolCircuit.TreeCircuit.depth, BoolCircuit.TreeCircuit.depth_le_maxDepth, BoolCircuit.TreeCircuit.depth_node, BoolCircuit.TreeCircuit.ind, BoolCircuit.TreeCircuit.maxDepth, BoolCircuit.TreeCircuit.maxFanin, BoolCircuit.TreeCircuit.maxFaninL, BoolCircuit.TreeCircuit.maxFanin_le_maxFaninL, BoolCircuit.TreeCircuit.maxFanin_node, BoolCircuit.TreeCircuit.size, BoolCircuit.TreeCircuit.size_node, BoolCircuit.TreeCircuit.sumSize_add_length_le, BoolCircuit.TreeCircuit.sumSize_nil}
\difficulty{5}
For every Boolean circuit $C$ on $n$ variables whose maximum fanin is at most $2$, the
size of $C$ plus $1$ is at most $2$ raised to the power of its depth plus $1$.
\end{theorem}

%%% added by the blueprint pipeline
\subsection*{Circuit Complexity --- Foundations: Further results}

\begin{definition}[Disjunction circuits]
\lean{BoolCircuit.NOrCircuit}
\label{BoolCircuit.NOrCircuit}
\leanok
A disjunction circuit on $n$ variables is either a clause consisting of a list of
literals whose variable indices are pairwise distinct, or a node consisting of a list of
conjunction circuits on the same $n$ variables.
\end{definition}
